\documentclass[11pt,twoside]{article}
\newcommand{\thang}[1]{{\color{blue}\small thang:#1}}

\usepackage{amssymb}
\usepackage{amsmath}
\usepackage{latexsym}
\usepackage{color}
\usepackage{graphics}
\usepackage{xspace}

% Commands for special characters
\newcommand{\mybackslash}{\char'134}
\newcommand{\myleftbracket}{\char'173}
\newcommand{\myrightbracket}{\char'175}
\newcommand{\mypercent}{\char'045}
\newcommand{\myunderscore}{\char'137}

% The 'ifthen' package supports Boolean flags
\usepackage{ifthen}
% The 'solutions' flag determines whether this is the original handout
%    or a solution
\newboolean{solutions}
\setboolean{solutions}{True}  % Default is original handout

% Uncomment the next line before starting on the solutions
\setboolean{solutions}{True}

\newcommand{\coursenumber}{CS 4104}
\newcommand{\coursetitle}{Data and Algorithm Analysis}
\newcommand{\docdate}{August 30, 2023}
\newcommand{\duedate}{September 11, 2023}
\newcommand{\homeworknumber}{1}

% The document title depends on whether these are solutions
\ifthenelse{\boolean{solutions}}{% Then
	\newcommand{\doctitle}{Solutions to Homework Assignment \homeworknumber}
}{% Else
	\newcommand{\doctitle}{Homework Assignment \homeworknumber}
}

% The document date depends on whether these are solutions
\ifthenelse{\boolean{solutions}}{% Then
	\renewcommand{\docdate}{\duedate}
}{% Else
}

% If you are the author, so put your name here
\renewcommand{\author}{An Author}

\pagestyle{myheadings}
\markboth{\hfill\doctitle\hfill\docdate}{\docdate\hfill\doctitle\hfill}

\addtolength{\textwidth}{1.00in}
\addtolength{\textheight}{1.00in}
\addtolength{\evensidemargin}{-1.00in}
\addtolength{\oddsidemargin}{-0.00in}
\addtolength{\topmargin}{-.50in}
\setlength{\footskip}{0pt}

\newcommand{\polyreduce}{\leq_{\mathrm{P}}}

\newcounter{problem}
\setcounter{problem}{0}
\newcommand{\problem}[1]{%
	\refstepcounter{problem}\noindent\textbf{[#1] \arabic{problem}.}}

\newcommand{\solution}{\bigskip\hrule\bigskip}
\newcommand{\problembreak}{\bigskip\hrule\bigskip}

\renewcommand{\theenumi}{\textbf{\Alph{enumi}}}
\renewcommand{\theenumii}{\textbf{\roman{enumii}}}

\newcommand{\emptystring}{\lambda}

% Pseudocode comment symbol
\newcommand{\comment}{\textbf{//}\ \ }

% Pseudocode line numbering
\newcounter{pseudocode}
\newcommand{\firstline}{\setcounter{pseudocode}{0}\linenumber}
\newcommand{\linenumber}{\refstepcounter{pseudocode}\thepseudocode}
\newcommand{\pseudoindent}{\hspace*{26pt}}

\newcommand{\nil}{\mbox{\textsc{nil}}}

% Mathematical symbols
\newcommand{\grid}{\mathcal{G}}  % Grid graph
\newcommand{\integers}{\mathbb{Z}}  % Integers
\newcommand{\naturals}{\mathbb{N}}  % Natural numbers
\newcommand{\reals}{\mathbb{R}}  % Real numbers

% Probability
\newcommand{\expect}[1]{\mathbf{E}\left[#1\right]}
\newcommand{\prob}[1]{\mathrm{Pr}\left[#1\right]}
\newcommand{\var}[1]{\mathrm{Var}\left[#1\right]}

% Logic
\newcommand{\NOT}[1]{\neg #1}
\newcommand{\AND}{\wedge}
\newcommand{\OR}{\vee}
\newcommand{\clause}[1]{\left(#1\right)}

\newlength{\problemoffset}
\setlength{\problemoffset}{0.5in}

% Decision problem macro
% A command for formatting decision problems a la Garey and Johnson
\newcommand{\decision}[3]{%     \decision{NAME}{INSTANCE}{QUESTION}
	\begin{list}{}{
			\setlength{\leftmargin}{\problemoffset}
			\setlength{\rightmargin}{\problemoffset}
			\setlength{\parsep}{0pt}
			\setlength{\itemsep}{2pt}
			\setlength{\topsep}{\itemsep}
			\setlength{\partopsep}{\itemsep}
		}
		\item
		{\textsc{#1}}
		\item
		{INSTANCE: #2}
		\item
		{QUESTION: #3}
	\end{list}
}

% Optimization problem macro
\newcommand{\optimization}[3]{%  \optimization{NAME}{INSTANCE}{SOLUTION}
	\begin{list}{}{
			\setlength{\leftmargin}{\problemoffset}
			\setlength{\rightmargin}{\problemoffset}
			\setlength{\parsep}{0pt}
			\setlength{\itemsep}{2pt}
			\setlength{\topsep}{\itemsep}
			\setlength{\partopsep}{\itemsep}
		}
		\item
		{\rule{0pt}{14pt}\textsc{#1}}
		\item
		{INSTANCE: #2}
		\item
		{SOLUTION: #3}
	\end{list}
}

\newcommand{\reaches}{\leadsto}

% Table layout
\newcommand{\toprule}{\rule[11pt]{0pt}{2pt}}
\newcommand{\bottomrule}{\rule[-6pt]{0pt}{0pt}}
\newenvironment{raggedpars}[1][2.0in]{%
	\begin{minipage}[t]{#1}\raggedright\toprule}%
	{\bottomrule\end{minipage}}

% Dynamic programming macros
\newlength{\arrowwidth}
\setbox3=\hbox{$\nwarrow$}
\setlength{\arrowwidth}{\wd3}
\newcommand{\optnwarrow}[1]{\if1#1\nwarrow%
	\else\rule{\arrowwidth}{0pt}\fi}
\newcommand{\optuparrow}[1]{\if1#1\uparrow%
	\else\rule{\arrowwidth}{0pt}\fi}
\newcommand{\optleftarrow}[1]{\if1#1\leftarrow%
	\else\rule{\arrowwidth}{0pt}\fi}
% Use \tablebox to specify any combination of arrows, plus the value
\newcommand{\tablebox}[4]{%
	\setlength{\arraycolsep}{0.0pt}%
	\begin{array}{cc}%
		\optnwarrow{#1} & \optuparrow{#2} \\%
		\optleftarrow{#3} & #4%
	\end{array}%
}
% Use \tableboxred to specify any combination of arrows, plus the value
% The value will be red; arrows are made red if 2 used instead of 1
\newcommand{\optnwarrowred}[1]{\if1#1\nwarrow%
	\else\if2#1\textcolor{red}{\nwarrow}\else\rule{\arrowwidth}{0pt}\fi\fi}
\newcommand{\optuparrowred}[1]{\if1#1\uparrow%
	\else\if2#1\textcolor{red}{\uparrow}\else\rule{\arrowwidth}{0pt}\fi\fi}
\newcommand{\optleftarrowred}[1]{\if1#1\leftarrow%
	\else\if2#1\textcolor{red}{\leftarrow}\else\rule{\arrowwidth}{0pt}\fi\fi}
\newcommand{\tableboxred}[4]{%
	\setlength{\arraycolsep}{0.0pt}%
	\begin{array}{cc}%
		\optnwarrowred{#1} &%
		\optuparrowred{#2} \\%
		\optleftarrowred{#3} &%
		\textcolor{red}{#4}%
	\end{array}%
}

% Allow really sloppy paragraphs that do not generate overfull or
%    underfull hbox's
\newenvironment{SLOPPY}{\begin{sloppypar}\hbadness=10000}{\end{sloppypar}}

% Definitions for this homework
\newcommand{\extent}[1]{\mathrm{extent}(#1)}
\newcommand{\opt}[2]{\mbox{\textsc{opt}}(#1,#2)}
\newcommand{\gap}{\mbox{\texttt{-}}}
\newcommand{\rewriterule}[2]{#1\to #2}
\newcommand{\rewrites}[2][]{\mathop{\Longrightarrow}\limits_{#1}^{#2}}
\newcommand{\reduces}{\leq}
\newcommand{\polyreduces}{\leq_P}
\newcommand{\mathsc}[1]{\mbox{\normalfont\textsc{#1}}}
\newcommand{\NP}{\mathcal{NP}}
\renewcommand{\P}{\mathcal{P}}

\begin{document}
	
	\thispagestyle{empty}
	
	\begin{center}
		\begin{tabular}{lcr}
			\multicolumn{3}{c}{\Large\textbf{\coursenumber}}
			\\[0.04in]
			\multicolumn{3}{c}{\Large\textbf{\doctitle}}
			% If these are solutions, then include the author's (student's) name
			\ifthenelse{\boolean{solutions}}{% Then
				\\[0.04in]\multicolumn{3}{c}{\large\textbf{\author}}
			}{} % Else omit
			% If these are solutions, then the date is the due date
			\ifthenelse{\boolean{solutions}}{% Then
				\\[0.10in]\multicolumn{3}{c}{\duedate}
			}{% Else, put given and due dates
				\\[0.10in]
				\textbf{Given:} \docdate
				& \hspace*{1.0in} &
				\textbf{Due:} \duedate
			}
		\end{tabular}
	\end{center}
	
	% If these are solutions, then we do not include the directions
	\ifthenelse{\boolean{solutions}}{} % No directions
	{
		% Original document includes directions
		\begingroup % This allows an argument that contains multiple paragraphs
		\paragraph{General directions.}
		The point value of each problem is shown in [ ].
		Each solution must include all details and
		an explanation of why the given solution is correct.
		\textbf{\textcolor{red}{In particular,
				write complete sentences.
				A correct answer without an explanation is worth no credit.}}
		The completed assignment must be submitted on Canvas as a PDF by 5:00 PM
		on \duedate.
		\textbf{No late homework will be accepted.}
		
		\paragraph{Digital preparation of your solutions is mandatory.}
		Use of \LaTeX\ is required.
		Also,
		\textbf{please include your name.}
		
		\paragraph{Use of \LaTeX\ (required).}
		\begin{itemize}
			\item Retrieve this \LaTeX\ source file,
			named
			\texttt{homework\homeworknumber.tex},
			from the course web site.
			\item Rename the file
			\textit{$<$Your VT PID$>$}\texttt{{\myunderscore}solvehw\homeworknumber.tex},
			For example,
			for the instructor,
			the file name would be
			\texttt{hoang{\myunderscore}solvehw\homeworknumber.tex}.
			
			\item
			Use a \textbf{text editor}
			(such as \texttt{vi}, \texttt{emacs}, or \texttt{pico})
			to accomplish the next three steps.
			Alternately,
			use Overleaf as your \LaTeX\ platform.
			
			\item
			Uncomment the line
			
			\texttt{{\mypercent} 
				{\mybackslash}setboolean{\myleftbracket}solutions{\myrightbracket}%
				{\myleftbracket}True{\myrightbracket}}
			
			in the document preamble by deleting the \texttt{\mypercent}.
			
			\item
			Find the line
			
			\texttt{{\mybackslash}renewcommand%
				{\myleftbracket}{\mybackslash}author{\myrightbracket}%
				{\myleftbracket}Thang Hoang{\myrightbracket}}
			
			and replace the instructor's name with your name.
			
			\item
			Enter your solutions where you find
			the \LaTeX\ comments
			
			\texttt{{\mypercent} PUT YOUR SOLUTION HERE}
			
			\item
			Generate a PDF and turn it in on Canvas by 5:00 PM on \duedate.
		\end{itemize}
		\endgroup
		
		\problembreak
		
		\newpage
		
	}
	
	\bigskip
	
	
	\problem{30}
	Find an algorithmic growth rate (and justify your answer in complete sentences) for the following two parts.
	
	{\bfseries
		\begin{enumerate}
			
			\item
			Find a growth rate  such that the time complexity
			is squared whenever the instance size is doubled. In other words, if $T(n) = X$, then $T(2n)= X^2$. \textbf{(15 pt)}
			
			\item
			Find a growth rate  such that the time complexity
			is cubed whenever the instance size is doubled. In other words, if $T(n) = X$, then $T(2n)= X^3$. \textbf{(15 pt)}
			
			
			
		\end{enumerate}
	}
	
	\solution
	% PUT YOUR SOLUTION HERE
	
	\begin{enumerate}
		
		\item
		$C^n$ for any constant $C$.
		Because $C^{2n} = C^nC^n = (C^n)^2$.
		
		\item
		$C^{n^{\log_2 3}}$ OR $C^{3^{\log_2 n}}$ for any constant $C$.
		
		A useful approach to tackling a problem like this is to start by trying to bracket the solution.
		From part (a), we know that $C^n$ is too small.
		$C^{2n}$ is still too small because $C^2$ is just another constant, which means that when you double $n$ you are just squaring the result.
		However, hopefully it is a short step to recognizing that $C^{(n^2)}$ is too big, since $C^{((2n)^2)} = (C^{(n^2)})^4$.
		So we have a good bracket. The next likely observation is that you are searching for an exponent that is a function of $n$ such that it triples
		when $n$ doubles. Why?
		Because if we have a function that triples when $n$ doubles, then we have an exponent that, when $C$ is taken to that exponent, gives a function that cubes when $n$ doubles, which is the answer to
		our problem. Since $(2n)^{\log_2 3} = 2^{\log_2 3} n^{\log_2 3}$, and $2^{\log_2 3} = 3$, we get $(2n)^{\log_2 3} = 3[n^{\log_2 3}$]. This is what we want.
		
	\end{enumerate}
	
	
	
	\problembreak
	
	\problem{30}
	Professor Alex uses the following algorithm for merging $k$ sorted lists, each having $n/k$ elements.
	She takes the first list and merges it with the second list using a linear-time algorithm for merging two sorted lists, such as the merging algorithm used in merge sort.
	Then, she merges the resulting list of $2n/k$ elements with the third list, merges the list of $3n/k$ elements that results with the fourth list, and so forth, until she ends up with a single sorted list of all elements.
	{\bfseries
		
		
		
		Analyze the worst-case time complexity of Professor Alex’s algorithm in terms of $n$ and $k$.
		Use complete sentences to explain your reasoning.
		Your final result should be a $\Theta$ expression.
		
		
	}
	
	\solution
	
	% PUT YOUR SOLUTION HERE
	
	Merging the first two lists, each of $n/k$ elements, takes $2n/k$ time. Merging the resulting $2n/k$ elements with the third list of $n/k$ elements takes $3n/k$ time and so on. Thus for a total of $k$ lists, we have:
	\begin{eqnarray*}
		T(n) & = &2n/k + 3n/k + \cdots + kn/k
		\\
		& = &\sum_{i=2}^{k} in/k
		\\
		& =  & n/k \sum_{i=2}^{k} i
		\\
		& = & n/k (k+2)(k-1)/2
		\\
		& = &
		\Theta(nk)
	\end{eqnarray*}
	
	
	
	
	\problembreak
	\problem{40}
	Consider the problem of computing  $a^n$, where $a > 0$ and $n\in\naturals$.
	{\bfseries
		\begin{enumerate}
			
			\begin{SLOPPY}
				
				\item
				State the computational problem {\normalfont\textsc{Exponentiation}}
				in formal Instance/Solution (or Input/Output) form
				as is done in the lecture notes. \textbf{(10 pt)}
				
			\end{SLOPPY}
			
			\item	Write CLRS pseudocode for a divide-and-conquer algorithm for the exponentiation problem of computing  $a^n$ where $a > 0$ and $n$ is a positive integer. \textbf{(10 pt)}
			\item	Define and solve a recurrence for the number of multiplications made by your algorithm. \textbf{(10 pt)}
			\item	How does your algorithm compare with a brute-force algorithm for this problem? \textbf{(10 pt)}
		\end{enumerate}
	}
	
	\solution
	\problembreak
	% PUT YOUR SOLUTION HERE
	\begin{enumerate}
		\item 
		The verbal description of \textsc{Exponentiation}
		translates directly into a formal problem,
		as follows:
		\optimization{Exponentiation}%
		{Positive integer $a$ and integer $n\in\naturals$.}%
		{The integer $a^n$.}
		
		\item
		The pseudocode can be found in Figure~\ref{figure:exponentiationdivideandconquer}.
		\begin{figure}[tp]
			\begin{tabbing}
				\textsc{Exponentiation}$(a, n)$ \\
				\firstline \qquad\=
				\comment Computes $a^n$ by a divide-and-conquer algorithm. 
				\\
				\linenumber\>
				\comment Input: A positive number a and a positive integer n. Output: The value of $a^n$
				\\[6pt]
				\linenumber\>
				\textbf{if} \= $n=0$
				\\
				\linenumber\>
				\> \textbf{return} $1$
				\\[6pt]
				\linenumber\>
				\textbf{if} \= $n=1$
				\\
				\linenumber\>
				\> \textbf{return} $a$
				\\[4pt]
				\linenumber\>
				\textbf{return} $\mbox{\textsc{Exponentiation}}(a, \lfloor n/2 \rfloor) \cdot {\textsc{Exponentiation}}(a, \lceil n/2 \rceil)$
			\end{tabbing}
			\caption{Pseudocode
				for divide-and-conquer algorithm
				for the \textsc{Exponentiation} problem.}
			\label{figure:exponentiationdivideandconquer}
		\end{figure}
		
		\item The recurrence for the number of multiplications is:
		\begin{eqnarray*}
			T(n) & = & T(\lfloor n/2 \rfloor) + T(\lceil n/2 \rceil) \text{ for } n > 1, M(1) = 0 , n = 2^k
			\\
			T(n) & = & 2 * T(2^k-1) + 1 
			\\
			T(n) & = & n-1
			\\
			T(n) & = & O(n)
		\end{eqnarray*}
		\item Though the algorithm makes the same number of multiplications as the brute-force method, it has to be considered inferior to the latter because of the recursion overhead.
	\end{enumerate}
	
\end{document}
